% Generated draft: requires mathematical and source review.
\documentclass[11pt]{article}
\usepackage[margin=25mm]{geometry}
\usepackage{fontspec}
\setmainfont{Latin Modern Roman}
\usepackage{amsmath,amssymb,mathtools,unicode-math}
\setmathfont{Latin Modern Math}
\usepackage{enumitem,needspace,xurl,hyperref}
\setlength{\parindent}{0pt}
\setlength{\parskip}{5pt}
\setlength{\emergencystretch}{3em}
\newcommand{\N}{\mathbb{N}}
\newcommand{\Z}{\mathbb{Z}}
\newcommand{\Q}{\mathbb{Q}}
\newcommand{\R}{\mathbb{R}}
\newcommand{\C}{\mathbb{C}}
\newcommand{\F}{\mathbb{F}}
\newcommand{\A}{\mathbb{A}}
\newcommand{\PP}{\mathbb{P}}
\newcommand{\E}{\mathbb{E}}
\newcommand{\eps}{\varepsilon}
\newcommand{\dd}{\,\mathrm{d}}
\begin{document}
\section*{UnsolvedMath ID: 30005809}
\textbf{Source problem code:} catalogue entry
\subsection{Definitions and mathematical statement}
For an integer \(n\geq 2\), consider indexed points \(p_1,\ldots,p_n\) in the closed unit square \([0,1]^2\subset\R^2\) and full affine lines \(\ell_1,\ldots,\ell_n\subset\R^2\) satisfying \(p_i\in\ell_i\) for every \(i\). Interpret distance as Euclidean distance:
\[
d(p,\ell)=\inf_{x\in\ell}\|p-x\|_2,
\qquad
\delta((p_i,\ell_i)_{i=1}^n)=\min_{\substack{1\leq i,j\leq n\\i\ne j}}d(p_i,\ell_j).
\]
The minimum is over ordered pairs of distinct indices. The supplied reconstruction imposes no explicit distinctness assumptions on the points or lines and no prohibition on additional incidences \(p_i\in\ell_j\) with \(i\ne j\).

Define the extremal quantity
\[
D_n=\sup\left\{\delta((p_i,\ell_i)_{i=1}^n):p_i\in[0,1]^2,\ \ell_i\text{ is an affine line in }\R^2,\ p_i\in\ell_i\text{ for all }i\right\}.
\]
The reconstructed problem asks to determine \(D_n\), and in particular to prove or disprove the assertion
\[
\forall\varepsilon>0\ \exists N_{\varepsilon}\in\N,\ N_{\varepsilon}\geq2,\ \forall n\geq N_{\varepsilon},\qquad D_n\leq n^{-1+\varepsilon}.
\]
Equivalently, the displayed bound must hold for every admissible configuration once \(n\) is sufficiently large depending only on \(\varepsilon\). The supremum expresses the largest possible separation without making an additional assertion about attainment.

This formulation follows the supplied reconstructed statement. The supplied original excerpt asks for the smallest point-to-line distance with unequal indices but does not itself specify the unit-square normalization, the paired incidences, or the asymptotic bound. The record labels the reconstruction as checked against the official report, while its supplied literature assessment says that normalization and hypotheses are missing. This tension remains unresolved; the formulation above is conditional on the reconstruction accurately representing the intended problem.
\subsection{Short English statement}
Among \(n\) points in the closed unit square, each paired with a line through it, maximize the minimum Euclidean distance from a point to a line with a different index. The reconstructed question asks whether this separation is at most \(n^{-1+\varepsilon}\) for every fixed \(\varepsilon>0\) and all sufficiently large \(n\). Its supplied open status and reconstruction require independent review.
\subsection{Sources}
\begin{itemize}
\item Supplied local UnsolvedMath registry record, permanent ID \(30005809\), title ``Point-to-Nonincident-Line Separation,'' source problem code \texttt{catalogue entry}. The supplied statement and clean statement agree. The record labels the statement \texttt{corrected\_verified} and says it was reconstructed and checked against the official report; this is a supplied verification claim, not an independent verification here.
\item \emph{Discrete Geometry} (2024), Oberwolfach Reports, pp.\ 137--202, DOI \texttt{10.4171/owr/2024/3}. Supplied source URL and literature-source link: \url{https://doi.org/10.4171/owr/2024/3}. The original excerpt associates the question with ``Lower bounds for incidences,'' Dmitrii Zakharov, joint work with Alex Cohen and Cosmin Pohoata. No separate proposer is specified.
\item Supplied literature assessment dated \texttt{2026-08-22}: status ``open.'' It states that the extracted distance question lacks normalization and hypotheses needed to identify a definite extremal theorem, describes the report as presenting an open incidence problem, and reports finding no authoritative resolution of the exact formulation. This is supplied dated curation, not an independently verified literature assessment. Its source list contains the report linked above.
\end{itemize}
\end{document}
